\section{A same-shell cascade obstruction}
\label{sec:fan-obstruction}

The first tempting route to \cref{eq:accelerated-revisit-target} is to charge
one gate batch to one new root-distance shell.  That statement is false on the
smallest nontrivial cyclic example.

\begin{proposition}[A tailed triangle activates twice in one shell]
\label{prop:tailed-triangle}
Let the vertices be $o,u,v,z$, with edges
\begin{equation}
 \{o,u\},\quad\{o,v\},\quad\{u,v\},\quad\{v,z\},
 \label{eq:tailed-triangle-edges}
\end{equation}
and seed $o$.  At
\begin{equation}
 \alpha=\frac3{20},
 \qquad
 \rho=\frac1{10},
 \label{eq:tailed-triangle-parameters}
\end{equation}
the exact boundary gate first admits $u$ and then admits $v$, although
\begin{equation}
 \operatorname{dist}(o,u)=\operatorname{dist}(o,v)=1.
 \label{eq:tailed-triangle-same-shell}
\end{equation}
Hence activation time is not bounded by root distance.
\end{proposition}

\begin{proof}
The degrees in the order $(o,u,v,z)$ are $(2,2,3,1)$.  Write
$g(x):=Qx-\ell_\rho$, so a boundary coordinate violates exactly when its
slack is negative.  On $U_0=\{o\}$, direct settlement gives
\begin{equation}
 x_o^0=\frac{24}{115\sqrt2}.
 \label{eq:tailed-triangle-first-value}
\end{equation}
Substitution into the two boundary slacks gives
\begin{equation}
 g_u(x^0)=-\frac{33\sqrt2}{4600}<0,
 \qquad
 g_v(x^0)=\frac{\sqrt3}{4600}>0.
 \label{eq:tailed-triangle-first-slacks}
\end{equation}
Thus $B_0=\{u\}$.  Settlement on $U_1=\{o,u\}$ gives
\begin{equation}
 x_o^1=\frac{334\sqrt2}{3045},
 \qquad
 x_u^1=\frac{44\sqrt2}{3045},
 \label{eq:tailed-triangle-second-values}
\end{equation}
and now
\begin{equation}
 g_v(x^1)=-\frac{21\sqrt3}{8120}<0.
 \label{eq:tailed-triangle-second-slack}
\end{equation}
Therefore $B_1=\{v\}$, proving the claim.
\end{proof}

The phenomenon is not bounded by a constant.  Define the tailed fan
$\mathcal F_m$ on
\begin{equation}
 V_m:=\{o,v_1,\ldots,v_m,z\}
 \label{eq:fan-vertices}
\end{equation}
by joining $o$ to every $v_i$, joining $v_i$ to $v_{i+1}$ for
$1\leq i<m$, and joining $v_m$ to $z$.  Thus
\begin{equation}
 d_o=m,\qquad d_{v_1}=2,\qquad
 d_{v_i}=3\ (2\leq i\leq m),\qquad d_z=1.
 \label{eq:fan-degrees}
\end{equation}
Every $v_i$ lies in the first root-distance shell.  The pendant $z$ removes
the second degree-two endpoint that would otherwise start a symmetric cascade.

\begin{theorem}[Arbitrarily long same-shell fan cascade]
\label{thm:fan-cascade}
Fix $\alpha\in(0,1)$.  For every integer $L\geq1$, there exist
$m>L$ and $\rho>0$ such that the exact boundary gate on $\mathcal F_m$ has
the first $L$ nonempty batches
\begin{equation}
 B_0=\{v_1\},\quad B_1=\{v_2\},\quad\ldots,\quad
 B_{L-1}=\{v_L\}.
 \label{eq:fan-singleton-batches}
\end{equation}
Every intermediate region has elimination width at most two, the entire graph
has root radius two, and the realized revisit factor satisfies
\begin{equation}
\mathfrak R_{\rm set}>\frac{L+1}{5}.
 \label{eq:fan-revisit-lower}
\end{equation}
Consequently, no bound on the literal gate's revisit factor can depend on
$\alpha$ alone.
\end{theorem}

\begin{proof}
Use the scaling
\begin{equation}
 \rho=\frac{\vartheta}{m},
 \qquad
 \vartheta_0:=\frac{1-\alpha}{4+2\alpha},
 \label{eq:fan-critical-scaling}
\end{equation}
and consider a fixed prefix length while $m\to\infty$.  Write
$x_o=X/\sqrt m$ and $x_{v_i}=Y_i/m$.  The settled root equation converges to
\begin{equation}
 X=\frac{2\alpha(1-\vartheta)}{1+\alpha}.
 \label{eq:fan-limit-root}
\end{equation}
For a root neighbor of degree $d$, the scaled shifted demand after the root
contribution is
\begin{equation}
 \psi_d(\vartheta)
 :=\alpha\left[
  \frac{(1-\alpha)(1-\vartheta)}{(1+\alpha)\sqrt d}
  -\vartheta\sqrt d
 \right].
 \label{eq:fan-limit-demand}
\end{equation}
The critical value in \cref{eq:fan-critical-scaling} is chosen so that
\begin{equation}
 \psi_3(\vartheta_0)=0,
 \qquad
 \psi_2(\vartheta_0)
 =\frac{\alpha\vartheta_0}{\sqrt2}>0.
 \label{eq:fan-critical-demands}
\end{equation}

Suppose the active fan prefix is $v_1,\ldots,v_j$.  In the same limit, its
scaled coordinates solve a tridiagonal Stieltjes system $T_jY=p$, with
diagonal
\begin{equation}
 a:=\frac{1+\alpha}{2},
 \label{eq:fan-core-diagonal}
\end{equation}
first off-diagonal magnitude
\begin{equation}
 h_1:=\frac{1-\alpha}{2\sqrt6},
 \label{eq:fan-first-coupling}
\end{equation}
all later off-diagonal magnitudes
\begin{equation}
 h:=\frac{1-\alpha}{6},
 \label{eq:fan-core-coupling}
\end{equation}
and right-hand side
\begin{equation}
 p=\bigl(\psi_2(\vartheta),
          \psi_3(\vartheta),\ldots,
          \psi_3(\vartheta)\bigr)^T.
 \label{eq:fan-core-load}
\end{equation}
At $\vartheta=\vartheta_0$, this is a positive multiple of $e_1$.
The irreducible Stieltjes inverse $T_j^{-1}$ is strictly positive, so the
last coordinate $Y_j^{(j)}$ is positive for every fixed $j$.

Here $\psi_d$ is the scaled violation demand
$m(\ell_\rho-Qx)$ before contributions from active path neighbors.  At the
critical value, a remote inactive degree-three vertex therefore has limiting
violation demand zero, whereas the next path vertex has limiting violation
demand
\begin{equation}
 \begin{cases}
  h_1Y_1^{(1)},&j=1,\\
  hY_j^{(j)},&j\geq2,
 \end{cases}
 \label{eq:fan-next-margin}
\end{equation}
which is strictly positive.  There are only finitely many margins for
$j<L$.  By continuity, one may choose
$\vartheta>\vartheta_0$ sufficiently close that
\begin{equation}
 \psi_3(\vartheta)<0,
 \qquad
 \psi_3(\vartheta)+h_jY_j^{(j)}(\vartheta)>0
 \quad(1\leq j<L),
 \label{eq:fan-separated-signs}
\end{equation}
where $h_j=h_1$ for $j=1$ and $h_j=h$ otherwise; also
$\psi_2(\vartheta)>0$.  Thus the degree-two endpoint wants to enter, every
remote degree-three coordinate does not, and precisely the next path vertex
wants to enter.

For every fixed $j$, the finite-$m$ scaled restricted system converges
entrywise to the displayed root equation and $T_j$ system.  All inequalities
in \cref{eq:fan-separated-signs} are strict, so one sufficiently large
$m>L$ preserves them simultaneously for $j=0,\ldots,L-1$.  Induction then
gives \cref{eq:fan-singleton-batches}.

Eliminating $z,v_m,v_{m-1},\ldots,v_1$ leaves at most the root and one path
neighbor after every pivot, so the fan and every active prefix have width at
most two.  All $v_i$ have root distance one and $z$ has distance two.
Finally, every region $U_0,\ldots,U_L$ contains the root and hence has
\begin{equation*}
 \operatorname{cvol}(U_j)\geq m+1.
\end{equation*}
The final support is contained in $V_m$, while
\begin{equation*}
 \operatorname{cvol}(V_m)
 =|V_m|+2|E_m|=(m+2)+4m=5m+2.
\end{equation*}
Therefore
\begin{equation*}
 \mathfrak R_{\rm set}
 \geq\frac{(L+1)(m+1)}{5m+2}
 >\frac{L+1}{5}.
\end{equation*}
\end{proof}

\begin{corollary}[A logarithmic dynamic-range loss is intrinsic]
\label{cor:fan-logarithmic-loss}
For every fixed $\alpha\in(0,1)$, the construction in
\cref{thm:fan-cascade} can be chosen so that
\begin{equation}
 \mathfrak R_{\rm set}
 =\Omega_\alpha\!\left(\log\frac1\rho\right).
 \label{eq:fan-log-revisit}
\end{equation}
Hence a fresh-refactor-and-rescan implementation of the exact gate can require
\begin{equation}
 \Omega_\alpha\!\left(
  \frac1\rho\log\frac1\rho
 \right)
 \label{eq:fan-log-work}
\end{equation}
degree work even on width-two, radius-two instances.  The statement does not
refute a $\widetilde O(1/(\rho\sqrt\alpha))$ theorem; it proves that a hidden
dynamic-range logarithm or an incremental settlement interface is necessary.
\end{corollary}

\begin{proof}
At $\vartheta_0$, forward elimination of
$T_jY=\psi_2(\vartheta_0)e_1$ has positive pivots at most $a$.
Consequently the endpoint margin in
\cref{eq:fan-next-margin} is bounded below by
\begin{equation}
 C_\alpha\left(\frac h a\right)^j
 \label{eq:fan-margin-decay}
\end{equation}
for a constant $C_\alpha>0$.  Choose
$\vartheta-\vartheta_0$ as a sufficiently small constant multiple of the
minimum margin through $L$.  The $T_j$ are uniformly strictly diagonally
dominant, so their inverse norms and the $\vartheta$-derivatives of all the
displayed demands are bounded in terms of $\alpha$ alone.  The finite-$m$
scaled root feedback and coefficient perturbation for a prefix of length at
most $L$ are $O_\alpha(L/m)$.  It is therefore sufficient to take
\begin{equation}
 m=O_\alpha\!\left(
  L\left(\frac a h\right)^L
 \right).
 \label{eq:fan-size-choice}
\end{equation}
Since $\rho=\vartheta/m$ and $\vartheta$ remains bounded away from zero,
\cref{eq:fan-size-choice} gives
$\log(1/\rho)=O_\alpha(L)$.  Combine this with
\cref{eq:fan-revisit-lower}.

For the work statement, each fresh stage scans or refactors a region
containing the degree-$m$ root and therefore costs $\Omega(m)$.  The chosen
$\rho$ satisfies $m=\Theta_\alpha(1/\rho)$ up to the harmless variation of
$\vartheta$, so the first $L$ stages give \cref{eq:fan-log-work}.
\end{proof}

The fan separates three statements that should not be conflated.  Support
radius can be constant while gate time is unbounded; constant elimination
width does not help if the same region is freshly settled after every event;
and the polylogarithmic target remains viable precisely because the cascade
margin decays geometrically.  This points away from a log-free batch theorem
and toward either a margin-aware amortization or retained Schur state.
